\documentclass[a4paper,11pt]{article}

% Pacotes básicos
\usepackage[utf8]{inputenc}
\usepackage[T1]{fontenc}
\usepackage[english]{babel}
\usepackage{amsmath, amssymb, amsthm}
\usepackage{geometry}
\usepackage{graphicx}
\usepackage{hyperref}

\usepackage{mathrsfs}
\usepackage{enumitem}

% Configurações de página
\geometry{
	a4paper,
	left=30mm,
	right=30mm,
	top=25mm,
	bottom=25mm,
}

\title{\textbf{Uma Estrutura Combinatória Baseada em Simetrias para a
		Conjectura de Goldbach}\\
	\large (Preprint)}
\author{
	Tiago Bandeira\thanks{tiagobandeira.goldbach2025@gmail.com}
}
\date{\today}

% Teoremas e definições
\newtheorem{theorem}{Teorema}[section]
\newtheorem{lemma}[theorem]{Lema}
\newtheorem{proposition}[theorem]{Proposição}
\newtheorem{corollary}[theorem]{Corolário}
\theoremstyle{definition}
\newtheorem{definition}[theorem]{Definição}
\theoremstyle{remark}
\newtheorem{remark}[theorem]{Observação}
\newtheorem{example}[theorem]{Exemplo}
\newtheorem{conjecture}{Conjectura}

\begin{document}
	
	\maketitle
	
	\begin{otherlanguage*}{english}
		\begin{abstract}
			Este trabalho apresenta uma estrutura combinatória para o estudo da
			Conjectura de Goldbach, baseada em simetrias entre números ímpares e em
			progressões aritméticas associadas ao Teorema de Dirichlet. Por meio da
			definição de uma sequência \(S_n(p)\), associada a cada primo \(p > 2\), e
			dos subconjuntos \(A_p = \{p + q \mid q \in Q_p\}\), onde \(Q_p\) é o
			conjunto dos primos em \(S_n(p)\), organiza-se o conjunto dos números pares
			maiores que dois em subconjuntos estruturados por simetria e progressão
			aritmética. Demonstra-se que: (i) \(S_n(p)\) contém infinitos primos
			(Teorema de Dirichlet); (ii) para todo par \(a > 6\), existe algum primo
			\(p < a/2\) com \(p \nmid a\) (condição necessária para \(a \in A_p\)); e
			(iii) a estrutura de propagação por deslocamento simétrico é bem definida.
			Identifica-se e expõe-se explicitamente a lacuna central que separa esses
			resultados de uma demonstração de Goldbach: garantir que \(a - p\) seja
			primo para algum \(p\) adequado é precisamente o enunciado da conjectura,
			e nenhum dos argumentos apresentados fecha esse passo. O trabalho é
			reposicionado, portanto, como uma contribuição estrutural que organiza o
			problema de forma alternativa e abre caminhos para investigações futuras
			por densidade analítica ou métodos de crivo.
		\end{abstract}
	\end{otherlanguage*}
	
	\begin{abstract}
		This paper presents a combinatorial framework for the study of the Goldbach
		Conjecture, based on symmetries among odd numbers and arithmetic progressions
		associated with Dirichlet's theorem. Through the definition of a sequence
		\(S_n(p)\) for each prime \(p > 2\), and subsets \(A_p = \{p + q \mid q
		\in Q_p\}\), where \(Q_p\) is the set of primes in \(S_n(p)\), the set of
		even numbers greater than two is organized into subsets structured by
		symmetry and arithmetic progression. It is proved that: (i) \(S_n(p)\)
		contains infinitely many primes (Dirichlet's theorem); (ii) for every even
		\(a > 6\), there exists a prime \(p < a/2\) with \(p \nmid a\) (a necessary
		condition for \(a \in A_p\)); and (iii) the propagation structure via
		symmetric displacement is well defined. The central gap separating these
		results from a proof of Goldbach is explicitly identified and discussed:
		ensuring that \(a - p\) is prime for some suitable \(p\) is precisely the
		statement of the conjecture, and none of the arguments presented closes
		this step. The work is thus repositioned as a structural contribution that
		organizes the problem in an alternative way and opens paths for future
		investigation via analytic density arguments or sieve methods.
	\end{abstract}
	
	\textbf{Palavras-chave:} conjectura de Goldbach; números primos; teoria dos
	números; progressões aritméticas; estrutura combinatória.
	
	\begin{quote}
		\textbf{Nota do Autor.}
		
		Este trabalho apresenta uma investigação exploratória sobre estruturas combinatórias associadas à Conjectura de Goldbach.
		
		O objetivo é desenvolver propriedades estruturais, simetrias e relações aditivas relacionadas às partições por números primos, sem reivindicar uma demonstração completa da conjectura.
		
		Resultados rigorosos e hipóteses heurísticas são explicitamente distinguidos ao longo do texto.
		
		Comentários, críticas e sugestões matemáticas são bem-vindos.
	\end{quote}
	
	%%%%%%%%%%%%%%%%%%%%%%%%%%%%%%%%%%%%%%%
	\begin{otherlanguage*}{english}
		\section{Introdução}
		
		A Conjectura de Goldbach, proposta em 1742 por Christian Goldbach em uma
		correspondência com Leonhard Euler, é um dos problemas mais antigos e
		famosos ainda não resolvidos da Teoria dos Números. Seu enunciado é simples
		e elegante: \textbf{todo número inteiro par maior que dois pode ser expresso
			como a soma de dois números primos}. Apesar de resistir a séculos de
		tentativas de demonstração, a conjectura já foi verificada
		computacionalmente para números muito grandes (da ordem de
		\(4 \cdot 10^{18}\))~\cite{oliveira}, mas uma prova geral e formal permanece
		desconhecida.
		
		Diversas abordagens foram desenvolvidas ao longo do tempo para tentar
		compreender a estrutura que envolve a distribuição dos números primos e
		sua relação com os números pares. Dentre os trabalhos mais relevantes,
		destacam-se os estudos baseados em análise matemática, combinatória
		aditiva, estatística dos primos e o uso de progressões aritméticas
		associadas ao Teorema de Dirichlet, que garante a existência de infinitos
		primos em certas classes aritméticas.
		
		O trabalho de Ellman (2000)~\cite{Ellman2000} propõe uma abordagem
		combinatória com espírito similar ao deste artigo: organizar as somas entre
		primos ímpares para cobrir os números pares. A distinção técnica central
		está em como os dois trabalhos tratam a relação entre a infinitude de primos
		numa progressão e a cobertura de um par específico. O Teorema de Dirichlet
		garante que existem infinitos primos em $S_n(p)$, mas não garante que
		$a - p$ seja primo para um $a$ fixo — distinção que este trabalho identifica
		explicitamente na Observação~\ref{rem:lacuna} e que é desenvolvida na
		Seção~\ref{sec:discussao}.
		
		O presente trabalho propõe uma \emph{estrutura combinatória} baseada na
		sequência \(S_n(p)\), composta por números ímpares associados a cada primo
		\(p > 2$ por relações de simetria, e nos subconjuntos \(A_p\) formados pelas
		somas \(p + q\) com \(q \in Q_p\) (os primos de \(S_n(p)\)). O objetivo
		\textbf{não é apresentar uma demonstração definitiva} da Conjectura de
		Goldbach, mas sim desenvolver uma estrutura matemática alternativa para
		organizá-la, identificar o que pode ser provado dentro dela e expor
		com precisão onde se encontra a lacuna que impede uma demonstração completa.
		
		A motivação para essa abordagem surge da observação das simetrias existentes
		entre os números ímpares menores que um dado número par, e da possibilidade
		de decompor o conjunto dos pares em subconjuntos construídos por somas do
		tipo \(p + q\), onde \(q\) percorre a sequência \(S_n(p)\). Essa
		decomposição revela uma estrutura interna que organiza os pares segundo
		progressões aritméticas, oferecendo uma nova perspectiva sobre o problema.
		
		O presente artigo está organizado da seguinte forma: na seção
		\textbf{Fundamentação Teórica} são apresentadas as definições formais,
		proposições, corolários e os teoremas clássicos que sustentam a construção
		da proposta. Na seção \textbf{Estrutura Combinatória para a Conjectura de
			Goldbach} são desenvolvidos os resultados que a estrutura permite
		demonstrar, seguidos da identificação explícita da lacuna central do
		argumento. Na sequência, apresenta-se a \textbf{Discussão e Implicações},
		na qual são analisadas as implicações dessa estrutura e os caminhos
		que ela abre para pesquisas futuras. Por fim, a \textbf{Conclusão} resume
		os achados e sugere direções para investigações posteriores.
		
		\subsection*{A Conjectura de Goldbach}
		\begin{conjecture}[Conjectura de Goldbach]
			Todo número inteiro par maior que dois pode ser expresso como a soma de
			dois números primos.
			
			Formalmente, para todo \(n \in \mathbb{N}\) tal que \(n > 2\) e \(n\) é
			par, existem \(p, q \in \mathbb{P}\) (com \(\mathbb{P}\) sendo o conjunto
			dos números primos) tais que:
			\[
			n = p + q
			\]
		\end{conjecture}
		
		%%%%%%%%%%%%%%%%%%%%%%%%%%%%%%%%%%%%%%%
		\section{Fundamentação Teórica}
		
		\begin{proposition}\label{prop:simetria_impares}
			Seja $a \in \mathbb{N}$, com $a > 2$ e $a$ par. Considere a sequência
			ordenada dos números ímpares estritamente menores que $a$:
			\[
			S = \{1, 3, 5, \dotsc, a - 1\}
			\]
			Então, para todo $k \in S$, vale que:
			\[
			a = k + (a - k)
			\]
			e ambos $k$ e $a - k$ pertencem a $S$, sendo simétricos em relação aos
			termos centrais da sequência.
		\end{proposition}
		
		\begin{proof}
			Os elementos de $S$ são descritos pela progressão aritmética:
			\[
			S = \{2n - 1 \mid n = 1, 2, \dotsc, \tfrac{a}{2}\}
			\]
			O número total de elementos é $\tfrac{a}{2}$, inteiro pois $a$ é par.
			
			Seja $k = 2n - 1$ arbitrário. Seu par simétrico é:
			\[
			a - k = a - (2n - 1) = (a - 2n) + 1
			\]
			Como $(a - 2n)$ é par, $(a - 2n) + 1$ é ímpar. Além disso,
			$1 \leq a - k \leq a - 1$, logo $a - k \in S$. A soma verifica:
			\[
			k + (a - k) = (2n - 1) + (a - (2n - 1)) = a. \qedhere
			\]
		\end{proof}
		
		\subsection{Exemplos}
		
		Seja $a = 12$: $S = \{1, 3, 5, 7, 9, 11\}$, com pares simétricos $1+11=12$,
		$3+9=12$, $5+7=12$.
		
		Seja $a = 10$: $S = \{1, 3, 5, 7, 9\}$, com pares simétricos $1+9=10$,
		$3+7=10$, $5+5=10$ (o termo central $5$ emparelha-se consigo mesmo).
		
		Dessa forma, qualquer número par $a > 2$ pode ser decomposto como a soma de
		dois ímpares menores que $a$, simétricos em relação aos termos centrais da
		sequência.
		
		\begin{proposition}\label{prop:distancia_simetrica}
			Seja $a \in \mathbb{N}$, com $a > 2$ e $a$ par, e seja
			$S = \{1, 3, 5, \dotsc, a - 1\}$.
			
			Sejam $b, c \in S$ dois termos simétricos em relação ao centro da sequência,
			e seja $i$ a distância posicional de $b$ até o centro. Então:
			\[
			|c - b| =
			\begin{cases}
				4i, & \text{se } |S| = \frac{a}{2} \text{ é ímpar} \\
				4i + 2, & \text{se } |S| = \frac{a}{2} \text{ é par}
			\end{cases}
			\]
		\end{proposition}
		
		\begin{proof}
			A sequência $S$ possui $N = \frac{a}{2}$ termos. Sejam $b = 2n-1$ (posição
			$n$) e $c = 2(N-n+1)-1 = 2N-2n+1$ (posição $N-n+1$). Então:
			\[
			|c - b| = |2N - 4n + 2| = 2|N - 2n + 1|.
			\]
			\begin{itemize}
				\item Se $N$ é ímpar, o centro único está em posição $m = \frac{N+1}{2}$,
				logo $i = |n - m|$ e $|N - 2n + 1| = 2i$, dando $|c-b| = 4i$.
				\item Se $N$ é par, os dois centros estão em $m_1 = \frac{N}{2}$ e
				$m_2 = \frac{N}{2}+1$. Para $b$ à esquerda do centro,
				$i = m_1 - n$ e $|N - 2n + 1| = 2i+1$, dando $|c-b| = 4i+2$.
				\qedhere
			\end{itemize}
		\end{proof}
		
		\begin{corollary}\label{coro:par_simetrico}
			Nas condições da Proposição~\ref{prop:distancia_simetrica}, o termo
			simétrico de $b$ satisfaz:
			\[
			c =
			\begin{cases}
				b + 4i, & \text{se } |S| \text{ é ímpar} \\
				b + 4i + 2, & \text{se } |S| \text{ é par}
			\end{cases}
			\]
		\end{corollary}
		
		\begin{theorem}[Teorema de Dirichlet para Progressões Aritméticas]%
			\label{teo:dirichlet}
			Sejam $a$ e $d$ inteiros positivos com $\mathrm{mdc}(a,d) = 1$. Então a
			progressão aritmética $a,\, a+d,\, a+2d,\, \dotsc$ contém infinitos números
			primos.
		\end{theorem}
		
		A demonstração pode ser encontrada em~\cite{apostol}. Este teorema garante
		a existência de infinitos primos em certas progressões, mas \emph{não}
		garante que um elemento específico de uma progressão seja primo — distinção
		fundamental para os limites do presente trabalho.
		
		\begin{theorem}[Postulado de Bertrand]\label{teo:postulado_de_bertrand}
			Para todo inteiro $n > 1$, existe pelo menos um primo $p$ tal que
			$n < p < 2n$.
		\end{theorem}
		
		A demonstração pode ser encontrada em Galvin~\cite{Galvin2013}.
		
		\begin{definition}[Sequência $S_n(p)$]\label{def:snp}
			Seja $p > 2$ um primo fixado. A sequência $S_n(p)$, com $n \in \mathbb{N}$,
			é o conjunto dos números ímpares da forma:
			\[
			S_n(p) =
			\begin{cases}
				p + 4n \\
				p + 4n + 2
			\end{cases}
			\quad \text{para todo } n \in \mathbb{N}.
			\]
			Observa-se que $S_n(p)$ cobre todos os ímpares maiores que $p$, conforme
			o Corolário~\ref{coro:par_simetrico}.
		\end{definition}
		
		\begin{proposition}\label{prop:snp_infinitos_primos}
			A sequência $S_n(p)$ contém infinitos números primos.
		\end{proposition}
		
		\begin{proof}
			Os dois sub-casos correspondem às progressões $p + 4n$ e $p + 4n + 2$,
			ambas da forma $p + 2k$ com $\mathrm{mdc}(p, 2) = 1$. Pelo
			Teorema~\ref{teo:dirichlet}, cada uma contém infinitos primos.
		\end{proof}
		
		\begin{remark}
			A Definição~\ref{def:snp} está relacionada diretamente ao
			Corolário~\ref{coro:par_simetrico}: os elementos de $S_n(p)$ são
			precisamente os ímpares que, somados a $p$, geram um número par maior
			que $2p$ por deslocamento simétrico.
		\end{remark}
		
		\begin{remark}
			Para $q \in S_n(p)$: ou $q$ é primo ímpar ou $q$ é composto ímpar. A
			Proposição~\ref{prop:snp_infinitos_primos} garante que o primeiro caso
			ocorre infinitas vezes, mas não determina \emph{quais} elementos de
			$S_n(p)$ são primos.
		\end{remark}
		
		\begin{proposition}
			Seja $p > 2$ um primo e $n \in \mathbb{N}$. Então $S_n(p) + p$ é par.
		\end{proposition}
		
		\begin{proof}
			$S_n(p)$ é ímpar e $p > 2$ é ímpar; a soma de dois ímpares é par.
		\end{proof}
		
		\begin{proposition}[Propriedade da divisibilidade na soma]%
			\label{prop:divisibilidade_na_soma}
			Sejam $a, b, c \in \mathbb{N}$ com $a \neq 0$. Então:
			\[
			a \mid (b + c) \text{ e } a \mid b \;\Longleftrightarrow\; a \mid c.
			\]
		\end{proposition}
		
		\begin{proof}
			($\Rightarrow$) Se $b + c = a \cdot f$ e $b = a \cdot g$, então
			$c = a(f-g)$, logo $a \mid c$.
			($\Leftarrow$) Se $b = ag$ e $c = ah$, então $b + c = a(g+h)$, logo
			$a \mid (b+c)$.
		\end{proof}
		
		%%%%%%%%%%%%%%%%%%%%%%%%%%%%%%%%%%%%%%%
		\section{Estrutura Combinatória para a Conjectura de Goldbach}
		
		\begin{proposition}[Simetria dos Múltiplos]\label{prop:simetria_dos_multiplos}
			Seja $a > 2$ par e $p$ primo com $p \mid a$. Então $S_n(p)$ é múltiplo
			de $p$ sempre que $a = p + S_n(p)$.
		\end{proposition}
		
		\begin{proof}
			Se $p \mid a$ e $a = p + q$, então pela
			Proposição~\ref{prop:divisibilidade_na_soma} temos $p \mid q$, ou seja,
			$q = S_n(p)$ é múltiplo de $p$.
		\end{proof}
		
		\begin{lemma}\label{lema:multip}
			Seja $p > 2$ primo e $q > p$ ímpar em $S_n(p)$. Se $q$ é primo ou
			$q$ não é múltiplo de $p$, então $2p \nmid (p + q)$.
		\end{lemma}
		
		\begin{proof}
			Se $q$ é primo e $q > p$, então $p \nmid q$. Logo, pela
			Proposição~\ref{prop:divisibilidade_na_soma}, $p \nmid (p+q)$, portanto
			$2p \nmid (p+q)$. O mesmo vale se $q$ é composto mas $p \nmid q$.
		\end{proof}
		
		\begin{corollary}
			Seja $p > 2$ primo. Se $q = S_n(p)$ é primo e $q > p$, então
			$p \nmid (p + S_n(p))$.
		\end{corollary}
		
		\begin{proof}
			Segue imediatamente do Lema~\ref{lema:multip}, pois $q > p$ exclui
			$q = p$.
		\end{proof}
		
		\begin{proposition}\label{prop:produto_p_menores_a_meio}
			Seja $a \in \mathbb{N}$ par com $a > 6$. Então:
			\[
			a \neq \prod_{p < a/2} p
			\]
			Em outras palavras, nenhum par $a > 6$ é igual ao produto de todos os
			primos estritamente menores que $a/2$.
		\end{proposition}
		
		\begin{proof}
			Para $a > 6$ temos $a/2 > 3$, logo o produto inclui ao menos $2 \cdot 3
			= 6$. Se $a/2 \geq 5$, o produto inclui $2 \cdot 3 \cdot 5 = 30 > a$
			(pois $a \geq 10$ e $30 > 10$). Para $a$ maior, o produto cresce
			multiplicativamente enquanto $a$ cresce linearmente, portanto a
			igualdade é impossível.
		\end{proof}
		
		\begin{proposition}[Existência de Primos Não Divisores]%
			\label{prop:existencia_de_primos_nao_divisores}
			Seja $a \in \mathbb{N}$ par com $a > 6$. Então existe pelo menos um
			primo $p < a/2$ tal que $p \nmid a$.
		\end{proposition}
		
		\begin{proof}
			Por contradição: se todo primo $p < a/2$ dividisse $a$, então $a$ seria
			múltiplo do produto de todos esses primos, contradizendo a
			Proposição~\ref{prop:produto_p_menores_a_meio}.
		\end{proof}
		
		\subsection{Definições dos Subconjuntos Estruturais}
		
		\begin{definition}[Conjunto $Q_p$]
			Seja $p > 2$ primo. Define-se:
			\[
			Q_p = \{q \in S_n(p) \mid q \text{ é primo}\}
			\]
			$Q_p$ é infinito pela Proposição~\ref{prop:snp_infinitos_primos}.
		\end{definition}
		
		\begin{definition}[Conjunto $A_p$]\label{def:Ap}
			Para $p > 2$ primo:
			\[
			A_p = \{p + q \mid q \in Q_p\}
			\]
			$A_p$ contém todos os pares gerados somando $p$ com os primos de $S_n(p)$.
		\end{definition}
		
		\begin{definition}[Conjunto $A$]
			\[
			A = \{a \in \mathbb{N} \mid a > 2 \text{ e } a \text{ é par}\}
			\]
		\end{definition}
		
		\subsection{Resultados Válidos sobre a Estrutura}
		
		\begin{proposition}[Condição Necessária para Cobertura]%
			\label{prop:condicao_necessaria}
			Seja $a > 6$ par. Então existe ao menos um primo $p < a/2$ com $p \nmid a$.
			Tal $p$ satisfaz a \emph{condição necessária} (mas não suficiente) para que
			$a \in A_p$, pois se $p \mid a$ e $a = p + q$, então $p \mid q$, de modo que
			$q$ não pode ser primo maior que $p$.
		\end{proposition}
		
		\begin{proof}
			A existência de tal $p$ segue diretamente da
			Proposição~\ref{prop:existencia_de_primos_nao_divisores}. A condição
			necessária segue da Proposição~\ref{prop:simetria_dos_multiplos}: se
			$p \mid a$, então $q = a - p$ é múltiplo de $p$, logo $q \notin Q_p$.
			Portanto, $a \in A_p$ exige $p \nmid a$.
		\end{proof}
		
		\begin{remark}\label{rem:lacuna}
			\textbf{(Identificação da Lacuna Central.)}
			A Proposição~\ref{prop:condicao_necessaria} mostra que a condição $p \nmid a$
			é necessária para $a \in A_p$, mas \emph{não} suficiente. Para concluir
			$a \in A_p$, é preciso garantir que $q = a - p$ seja primo. Isso é
			precisamente o enunciado da Conjectura de Goldbach — não decorre do Teorema
			de Dirichlet nem de qualquer resultado apresentado aqui. O Teorema de
			Dirichlet garante que existem infinitos primos em $S_n(p)$; \emph{não}
			garante que $a - p$ especificamente seja um deles.
		\end{remark}
		
		\begin{proposition}[Propagação Condicional por Deslocamento Simétrico]%
			\label{prop:propagacao}
			Seja $a = p + q \in A_p$ com $q \in Q_p$. Para todo $k \in \mathbb{N}$,
			se $q + 2k$ for primo (respectivamente $q - 2k > p$ e primo), então:
			\[
			a + 2k = p + (q + 2k) \in A_p \quad
			(\text{resp.}\ a - 2k = p + (q - 2k) \in A_p).
			\]
		\end{proposition}
		
		\begin{proof}
			Se $q + 2k$ é primo e $q + 2k > p$, então $q + 2k \in S_n(p)$ (pois
			$S_n(p)$ cobre todos os ímpares $> p$) e, sendo primo, $q + 2k \in Q_p$.
			Portanto $a + 2k = p + (q + 2k) \in A_p$. O caso $q - 2k$ é análogo.
		\end{proof}
		
		\begin{remark}
			A Proposição~\ref{prop:propagacao} é válida, mas condicional: ela afirma
			que \emph{se} existir um primo na posição deslocada, \emph{então} o par
			correspondente pertence a $A_p$. A existência desse primo não é garantida
			pelo argumento estrutural, e esta é exatamente a mesma lacuna identificada
			na Observação~\ref{rem:lacuna}.
		\end{remark}
		
		\begin{lemma}[Propagação Estrutural por Deslocamento Simétrico]
			Seja $a = p + q$ com $p \in \mathbb{P}$, $q \in Q_p$ e $q > p$. Se
			$q \pm 2k$ for primo (e $> p$), então $q \pm 2k \in Q_p$ e
			$a \pm 2k \in A_p$.
		\end{lemma}
		
		\begin{proof}
			Segue diretamente da Proposição~\ref{prop:propagacao} e da definição de
			$Q_p$.
		\end{proof}
		
		\subsection{A Lacuna Central e o Limite do Argumento}\label{sec:lacuna}
		
		Os resultados acima estabelecem:
		\begin{enumerate}
			\item Para todo $a > 6$ par, existe $p < a/2$ primo com $p \nmid a$
			(condição necessária para $a \in A_p$).
			\item $S_n(p)$ contém infinitos primos (Dirichlet), logo $Q_p$ é infinito
			e $A_p$ é infinito.
			\item Se $a \in A_p$ para algum $p$, então os pares $a \pm 2k$ pertencem
			a $A_p$ para todo $k$ tal que o deslocamento gere um primo.
		\end{enumerate}
		
		O que \textbf{não} foi demonstrado — e constitui a lacuna que impede que
		este trabalho seja uma prova de Goldbach — é:
		
		\begin{conjecture}[Cobertura Total — equivalente a Goldbach]\label{conj:cobertura}
			Para todo número par $a > 2$, existe ao menos um primo $p < a/2$ tal que
			$a \in A_p$, ou seja, tal que $a - p \in Q_p$ (i.e., $a - p$ é primo).
		\end{conjecture}
		
		\begin{remark}
			A Conjectura~\ref{conj:cobertura} é logicamente equivalente à Conjectura
			de Goldbach: dizer que existe $p$ primo com $a - p$ primo é exatamente
			dizer que $a$ é soma de dois primos. Portanto, qualquer argumento que
			pretenda demonstrá-la sem hipóteses adicionais está, em essência,
			tentando provar Goldbach diretamente.
		\end{remark}
		
		Os argumentos apresentados nas versões anteriores deste trabalho nos Lemas
		3.14 e 3.15 (Cobertura Estrutural e Contradição por Exclusão Total)
		continham um raciocínio circular: assumiam que a exclusão de $a$ de todos
		os $A_p$ contradiria Dirichlet por exigir que ``todos os deslocamentos
		deixem de gerar primos''. Mas Dirichlet garante apenas a infinitude de
		primos na progressão $S_n(p)$ — não que \emph{especificamente} $a - p$
		seja primo. A não pertença de $a$ a $A_p$ não requer que nenhum primo
		exista em $S_n(p)$; requer apenas que $a - p \notin \mathbb{P}$, o que é
		compatível com a infinitude de outros primos em $S_n(p)$.
		
		Analogamente, o Postulado de Bertrand garante um primo entre $a/2$ e $a$,
		mas não garante que ele seja da forma $a - p$ para algum primo $p$.
		
		\subsection{Resultados Condicionais}
		
		Enunciamos a seguir os resultados que \emph{seriam} válidos caso a
		Conjectura~\ref{conj:cobertura} fosse provada, a título de organização
		estrutural e motivação para investigações futuras.
		
		\begin{theorem}[Cobertura dos Números Pares — condicional]%
			\label{teo:cobertura_dos_pares}
			Assumindo a Conjectura~\ref{conj:cobertura}, seja
			\[
			B = \{(a_n)_{n \in \mathbb{N}} \mid a_n = p + q,\, p \in \mathbb{P},\,
			q \in Q_p\}.
			\]
			Então a união das sequências $a_n \in B$ cobre integralmente o conjunto
			dos números pares maiores que dois:
			\[
			\bigcup_{n \in \mathbb{N}} a_n = A.
			\]
			
			\emph{(Caso especial: $a = 4 = 2 + 2$. Para $p = 2$, $S_n(2)$ não satisfaz
				as condições de Dirichlet pois $\mathrm{mdc}(4,2) \neq 1$. O número $4$
				é tratado como caso particular via $a_1 = \{4\}$.)}
		\end{theorem}
		
		\begin{corollary}[Formulação Simétrica — condicional]\label{coro:par_simetrico_primos}
			Assumindo a Conjectura~\ref{conj:cobertura}: todo número par $a > 2$
			possui um par de primos $p$ e $q \in S_n(p)$ tal que $a = p + q$.
		\end{corollary}
		
		\begin{theorem}[Reformulação Simétrica da Conjectura de Goldbach]%
			\label{teo:goldbach_por_simetria_de_primos}
			As afirmações a seguir são equivalentes:
			\begin{enumerate}[label=(\roman*)]
				\item \emph{(Conjectura de Goldbach)} Todo par $a > 2$ é soma de dois
				primos.
				\item \emph{(Conjectura~\ref{conj:cobertura})} Para todo par $a > 2$,
				existe primo $p < a/2$ com $a - p \in Q_p$.
				\item \emph{(Cobertura)} $\bigcup_{p \in \mathbb{P}} A_p = A$.
			\end{enumerate}
		\end{theorem}
		
		\begin{proof}
			(i)$\Leftrightarrow$(ii): Seja $a = p + q$ com $p, q$ primos. Como $p$
			e $q$ são ambos maiores que $1$ e $p < a$, tomando $p < a/2$ (o menor
			dos dois), temos $q = a - p > a/2 \geq p$, logo $q > p$ e $q \in S_n(p)$,
			portanto $q \in Q_p$, dando (ii). A recíproca é imediata.
			
			(ii)$\Leftrightarrow$(iii): Por definição, $a \in A_p$ se e somente se
			$a - p \in Q_p$ para algum $p$.
		\end{proof}
		
		\begin{remark}
			O Teorema~\ref{teo:goldbach_por_simetria_de_primos} é uma reformulação,
			não uma prova. Ele mostra que a estrutura $A_p$ é uma linguagem alternativa
			para expressar Goldbach, o que pode ser útil para ataques futuros ao problema.
		\end{remark}
		
		%%%%%%%%%%%%%%%%%%%%%%%%%%%%%%%%%%%%%%%
		\section{Discussão e Implicações}\label{sec:discussao}
		
		\subsection{Interpretação Estrutural}
		
		A sequência $S_n(p)$ cobre todos os ímpares maiores que $p$, e $Q_p$ é o
		seu subconjunto de primos — infinito por Dirichlet. Os subconjuntos $A_p$
		organizam os pares maiores que $2p$ de acordo com o primo âncora $p$. A
		interseção $A_p \cap A_{p'}$ para diferentes primos $p, p'$ é não vazia em
		geral, o que reflete o fato de que um número par pode ser soma de mais de
		um par de primos (partições de Goldbach).
		
		A estrutura de propagação (Proposição~\ref{prop:propagacao}) é uma
		propriedade genuína dos conjuntos $A_p$: uma vez que $a \in A_p$, os pares
		adjacentes $a \pm 2$ pertencem a $A_p$ condicionalmente à primalidade de
		$q \pm 1$ — condição que está sujeita à distribuição dos primos e não pode
		ser garantida a priori.
		
		\subsection{Nota sobre Trabalhos Relacionados: Comparação com Ellman}
		\label{sec:ellman}
		
		O trabalho de Ellman~\cite{Ellman2000} explora uma ideia estruturalmente próxima
		à deste artigo: organizar as somas de primos ímpares em subconjuntos para cobrir
		os números pares. Vale registrar, com precisão técnica, o ponto em que os dois
		trabalhos divergem — não para criticar a iniciativa, mas porque a distinção é
		pedagogicamente útil para delimitar o escopo do presente artigo.
		
		Ellman argumenta que, como o Teorema de Dirichlet garante infinitos primos em
		progressões aritméticas e os subconjuntos de somas crescem indefinidamente, a
		cobertura total dos pares se seguiria. A limitação técnica desse argumento é que
		\emph{infinitos primos em $S_n(p)$} não implica que \emph{$a - p$ especificamente
			seja primo para um $a$ dado}. Em linguagem formal: Dirichlet garante $|Q_p| = \infty$,
		mas não garante $a - p \in Q_p$ para um $a$ fixo. A não pertença de $a$ a $A_p$
		é compatível com a existência de infinitos outros primos em $S_n(p)$.
		
		Esta distinção — entre \emph{existência assintótica} e \emph{cobertura pontual} —
		é precisamente o que a Observação~\ref{rem:lacuna} identifica como a lacuna central
		deste trabalho. Ao nomear essa distinção de forma explícita, este artigo posiciona
		a estrutura $S_n(p), A_p$ como uma contribuição organizacional honesta, sem
		reivindicar o que ela ainda não demonstra.
		
		\subsection{O Que Este Trabalho Demonstra}
		
		Os resultados plenamente demonstrados são:
		\begin{itemize}
			\item Para todo $a > 6$, existe $p < a/2$ com $p \nmid a$
			(Prop.~\ref{prop:existencia_de_primos_nao_divisores}) — condição
			necessária para $a \in A_p$.
			\item $Q_p$ é infinito para todo primo $p > 2$ (Dirichlet).
			\item A propagação dentro de $A_p$ funciona sempre que existem primos nas
			posições deslocadas (Prop.~\ref{prop:propagacao}).
			\item As três formulações do Teorema~\ref{teo:goldbach_por_simetria_de_primos}
			são logicamente equivalentes entre si e a Goldbach.
		\end{itemize}
		
		\subsection{O Que Não Foi Demonstrado}
		
		A afirmação central que permanece aberta é a
		Conjectura~\ref{conj:cobertura}: para todo $a$ específico, garantir que
		$a - p$ é primo para algum primo $p$. Este passo exige que o par $(p, a-p)$
		seja uma partição de Goldbach — que é exatamente o que a conjectura afirma.
		
		Nenhum argumento de infinitude (Dirichlet, Bertrand, densidade dos primos)
		é suficiente para garantir isso: eles estabelecem que existem muitos primos
		``próximos'' de $a/2$, mas não que algum deles seja complementar de um primo.
		
		\subsection{Caminhos para Investigações Futuras}
		
		A estrutura proposta sugere algumas direções promissoras:
		
		\begin{enumerate}
			\item \textbf{Argumento de densidade}: estimar a densidade de $A_p$ dentro
			de $A$ e mostrar que, para $a$ grande, a probabilidade de $a \notin
			\bigcup A_p$ tende a zero. Isso exigiria funções $L$ de Dirichlet e
			estimativas de crivo (Crivo de Selberg, por exemplo).
			
			\item \textbf{Análise das interseções}: investigar formalmente
			$|A_{p_1} \cap A_{p_2} \cap \dotsb \cap A_{p_k}|$ para coleções
			finitas de primos, conectando com as conjecturas de Hardy-Littlewood
			sobre pares de primos.
			
			\item \textbf{Verificação computacional}: os subconjuntos $A_p$ são
			gerados computacionalmente a partir de um primo $p$ e dos primos de
			$S_n(p)$. Uma análise empírica da cobertura progressiva $\bigcup_{p
				\leq P} A_p$ em função de $P$ e do alcance $a_{\max}$ pode revelar
			padrões para orientar o trabalho teórico.
			
			\item \textbf{Extensão a outros problemas}: a estrutura de simetria pode
			ser adaptada para a Conjectura de Goldbach fraca (já provada por
			Helfgott~\cite{Helfgott2013}) como teste de validação, ou para
			outros problemas da combinatória aditiva dos primos.
		\end{enumerate}
		
		%%%%%%%%%%%%%%%%%%%%%%%%%%%%%%%%%%%%%%%
		\section{Conclusão}
		
		Este trabalho desenvolveu uma estrutura combinatória alternativa para
		organizar a Conjectura de Goldbach, baseada na sequência \(S_n(p)\) e nos
		subconjuntos \(A_p = \{p + q \mid q \in Q_p\}\). Os resultados válidos
		obtidos são: (i) a existência garantida de um primo \(p < a/2\) com
		\(p \nmid a\) para todo par \(a > 6\), fornecendo a condição necessária
		para \(a \in A_p\); (ii) a infinitude de \(Q_p\) por Dirichlet; (iii) a
		propriedade de propagação condicional dentro de \(A_p\); e (iv) a
		equivalência lógica entre a Conjectura de Goldbach e a cobertura
		\(\bigcup A_p = A\).
		
		A lacuna central foi identificada e exposta com precisão: o passo de
		garantir que \(a - p\) seja primo para algum \(p\) adequado não decorre
		dos resultados anteriores e é logicamente equivalente à própria conjectura.
		Versões anteriores deste trabalho continham argumentos circulares nesse
		ponto (confusão entre ``existem infinitos primos em \(S_n(p)\)'' e
		``\(a - p\) é primo''), que foram removidos na presente versão.
		
		O trabalho é apresentado, portanto, não como uma demonstração de Goldbach,
		mas como uma contribuição estrutural que reapresenta o problema em uma
		linguagem de simetria e progressão aritmética, potencialmente útil como
		ponto de partida para abordagens por densidade analítica, crivos ou
		argumentos probabilísticos.
		
		%%%%%%%%%%%%%%%%%%%%%%%%%%%%%%%%%%%%%%%
		
		\section*{Agradecimentos}
		
		Ofereço este trabalho em homenagem a todos os meus professores de
		matemática, que contribuíram para minha formação e paixão pela matemática.
		
		Agradeço também às revisões críticas que motivaram a reformulação honesta
		do escopo deste trabalho.
		
		%%%%%%%%%%%%%%%%%%%%%%%%%%%%%%%%%%%%%%%
		
		\begin{thebibliography}{9}
			
			\bibitem{apostol}
			T. M. Apostol.
			\textit{Introduction to Analytic Number Theory}.
			Springer, 1976.
			
			\bibitem{paiva}
			C. D. C. Paiva et al.
			\textit{A busca pela prova da conjectura de Goldbach: explorando suas
				conquistas}.
			Caderno Pedagógico, 21(7): e5770--e5770, 2024.
			
			\bibitem{oliveira}
			T. Oliveira e Silva, S. Herzog, and S. Pardi.
			Empirical verification of the even Goldbach conjecture and computation of
			prime gaps up to \(4 \cdot 10^{18}\).
			\textit{Math.\ Comp.} \textbf{83} (2014), no.\ 288, 2033--2060.
			
			\bibitem{hefez}
			A. Hefez.
			\textit{Elementos de Aritmética}.
			Sociedade Brasileira de Matemática, 2006.
			
			\bibitem{Galvin2013}
			D. Galvin.
			\textit{Erdős's proof of Bertrand's postulate}, 2013.
			\url{https://www3.nd.edu/~dgalvin1/pdf/bertrand.pdf}
			
			\bibitem{Ellman2000}
			R. Ellman.
			\textit{A Proof of ``Goldbach's Conjecture''}, 2000.
			\url{https://arxiv.org/abs/math/0005185}
			
			\bibitem{Helfgott2013}
			H. A. Helfgott.
			\textit{Major arcs for Goldbach's theorem}, 2013.
			\url{https://arxiv.org/abs/1305.2897}
			
		\end{thebibliography}
		
	\end{otherlanguage*}
\end{document}